\section{Client Architecture}
\subsection{DAQ Nodes (STM32 Publishers)}
Each STM32 DAQ node operates as an MQTT publisher using the lwIP MQTT client (\texttt{lwip/apps/mqtt.h}) on top of the lwIP TCP/IP stack and FreeRTOS. The lwIP client was chosen over Paho Embedded C and FreeRTOS coreMQTT as it is integrated with the network stack already in use, adding no further dependency.

\paragraph{Publish Path} ~\\
Acquisition and publishing run in separate FreeRTOS tasks:
\begin{enumerate}[noitemsep]
  \item The ADCs sample continuously at approximately 30~kHz; DMA fills a ring buffer and signals the DAQ task when a block completes.
  \item The DAQ task scales and timestamps each block into a packet and places it on a ready queue. If no packet slot is free, the block is dropped rather than stalling acquisition.
  \item The MQTT task is the sole consumer of the ready queue. Every 5~s, it collects one 50~Hz cycle (20 samples at approximately 1~kHz) of each channel across consecutive packets, serialises one JSON burst per channel, and publishes the six messages at QoS~0.
  \item When built with \texttt{MQTT\_PUBLISH\_DAQ2}, the MQTT task also publishes every packet in binary on \texttt{daq2/<node\_id>/samples} (approximately 470 per second), at QoS~0. A packet that cannot be queued in the lwIP output buffer is dropped and counted rather than waited for.
  \item After each PTP servo update (on every Sync, about once a second), the MQTT task publishes the servo's state on \path{microgrid/<node_id>/ptp_status/servo}.
  \item Every packet is released back to the pool immediately after use, so acquisition never runs out of slots.
\end{enumerate}

\paragraph{Connection Management} ~\\
Each node uses a static IP address, and the broker address, set in \texttt{app/node\_config.h}. The MQTT task attempts a connection at most once per second while the Ethernet link is up, abandons an attempt that has not completed within 10~s, and reconnects automatically after broker loss. If the link goes down, the session is dropped. The keep-alive interval is 60~s. Telemetry is not buffered during disconnection; the resulting gap is visible in the historian, which is acceptable for a research testbed.

Client IDs follow the convention \texttt{mqttc-<node\_id>}. No credentials are configured, matching the broker's current anonymous access (Section~\ref{sec:security}).

\paragraph{Diagnostics} ~\\
Once per second, each node prints acquisition, MQTT, and Ethernet counters (connections, published and dropped messages, overruns, queue drops, invalid timestamps) and the PTP lock state to its debug UART.

\subsection{Historian Subscribers}
Two subscribers on the historian server divide the topic space between them: Telegraf subscribes to \texttt{microgrid/\#} (the JSON bursts and PTP status), and the \texttt{daq2-ingest} service to \texttt{daq2/+/samples} (the full-rate binary stream).

\paragraph{Telegraf} ~\\
Telegraf with the MQTT consumer input plugin is used as the bridge for the JSON topics. This provides a configuration-driven pathway from the broker to InfluxDB with no custom subscriber code. A regex processor extracts each topic segment into a tag (\texttt{node\_id}, \texttt{subsystem}, \texttt{metric}), and the payload's \texttt{ts} is used as the point timestamp. The as-built configuration (\texttt{telegraf/telegraf.conf}) is:

\begin{lstlisting}[language=toml]
[agent]
  flush_interval = "1s"
  precision = "1ns"

[[inputs.mqtt_consumer]]
  servers = ["tcp://mosquitto:1883"]
  topics = ["microgrid/#"]
  qos = 0
  data_format = "json"
  json_time_key = "ts"
  json_time_format = "unix_ns"
  topic_tag = "topic"

[[processors.regex]]
  order = 1
  [[processors.regex.tags]]
    key = "topic"
    pattern = "^microgrid/([^/]+)/([^/]+)/([^/]+)$"
    replacement = "${1}"
    result_key = "node_id"
  # ... likewise ${2} -> subsystem, ${3} -> metric

[[outputs.influxdb_v2]]
  urls = ["http://influxdb:8086"]
  token = "${INFLUXDB_TOKEN}"
  organization = "qut-microgrid"
  bucket = "telemetry"
\end{lstlisting}

Telegraf's JSON parser flattens the \texttt{v} array into fields \texttt{v\_0} to \texttt{v\_19}, so each burst is stored as a single InfluxDB point at the first sample's timestamp. A Starlark processor additionally splits each burst into 20 points, one per sample at \texttt{ts} $+ k \cdot$ \texttt{dt\_ns}. PTP status messages pass through unchanged as one point per Sync. The resulting schema is described in the InfluxDB design document.

\paragraph{daq2-ingest} ~\\
Telegraf cannot compute per-sample times from the binary packet header, so the full-rate stream has a dedicated subscriber: \texttt{daq2-ingest}, a Python service using the Eclipse Paho client. One receiver process subscribes to \texttt{daq2/+/samples} at QoS~0 and hands each node's packets to a worker process, which decodes them and writes them to InfluxDB. The service is described in the InfluxDB design document.
